% SPDX-License-Identifier: CC-BY-SA-4.0
% SPDX-FileCopyrightText: Copyright © 2026 Amanda Canizela Guimarães <amanda.canizela@gmail.com>
% SPDX-FileCopyrightText: Copyright © 2026 Ariel Inácio Jordão <arielijordao@gmail.com>
% SPDX-FileCopyrightText: Copyright © 2026 Lucca M. A. Pellegrini <lucca@verticordia.com>
% SPDX-FileCopyrightText: Copyright © 2026 Paulo Dimas Junior <paulo.junior.1478361@sga.pucminas.br>
% SPDX-FileCopyrightText: Copyright © 2026 Pedro Vitor Andrade <pedrovitor0826@gmail.com>

\documentclass[conference]{IEEEtran}
%Template version as of 6/27/2024

\usepackage{cite}
\usepackage{amsmath,amssymb,amsfonts}
\usepackage{algorithmic}
\usepackage{graphicx}
\usepackage{textcomp}
\usepackage{xcolor}

\usepackage{embedfile}
\usepackage[inline]{enumitem}
\usepackage{hyperref}
\usepackage{newtxtext,newtxmath}

\embedfile[
	desc={LaTeX source for this document},
	mimetype=text/plain
]{main.tex}

\embedfile[
	desc={BibTeX source file with references},
	mimetype=text/plain
]{references.bib}

\begin{document}

\title{A Quantitative Cache Evaluation of Select PolyBench Kernels}

\author{
    \IEEEauthorblockN{
        Amanda Canizela Guimarães\IEEEauthorrefmark{1},
        Ariel Inácio Jordão\IEEEauthorrefmark{2},
        Lucca Pellegrini\IEEEauthorrefmark{3},\\
        Paulo Dimas Junior\IEEEauthorrefmark{4},
        Pedro Vitor Andrade\IEEEauthorrefmark{5}
    }
    \IEEEauthorblockA{
        \textit{Instituto de Ciências Exatas e Informática}\\
        \textit{Pontifícia Universidade Católica de Minas Gerais}\\
        Belo Horizonte, Brazil \\
		\IEEEauthorrefmark{1}\href{mailto:amanda.canizela@gmail.com}{amanda.canizela@gmail.com}
		\IEEEauthorrefmark{2}\href{mailto:arielijordao@gmail.com}{arielijordao@gmail.com}
		\IEEEauthorrefmark{3}\href{mailto:lucca@verticordia.com}{lucca@verticordia.com}\\
		\IEEEauthorrefmark{4}\href{mailto:paulo.junior.1478361@sga.pucminas.br}{paulo.junior.1478361@sga.pucminas.br}
		\IEEEauthorrefmark{5}\href{mailto:pedrovitor0826@gmail.com}{pedrovitor0826@gmail.com}
    }
}

\maketitle

\begin{abstract}

	We present a systematic and reproducible evaluation of cache hierarchy
	sensitivities across five PolyBench kernels using the gem5 simulator. Our
	study separates realistic multi-level capacity presets from two orthogonal
	parameters  to isolate spatial and conflict effects. For each workload we
	execute a controlled sweep and collect a standard set of per-level miss
	rates and \textsc{mpki}, total miss latency, \textsc{ipc}/\textsc{cpi}.
	Three findings emerge. First, cache line size is the dominant factor:
	doubling the line often halves L1D miss rate and improves \textsc{ipc}
	where memory stalls are non-trivial. Second, associativity has minimal
	impact for these regular kernels, with the notable exception of
	direct-mapped caches. Third, capacity presets matter selectively when the
	hot working set approaches the \textsc{llc} boundary; otherwise spatial
	locality dominates. To strengthen reproducibility, we encode the full
	pipeline in a single build graph with pinned versions for all components,
	and statically link workloads. The resulting methodology and artifacts
	facilitate transparent and repeatable cache design-space exploration, and
	provide a baseline for future studies that consider larger datasets,
	alternative cache policies, or out-of-order cores.

\end{abstract}

\begin{IEEEkeywords}
    Memory hierarchy, cache memories, simulation, performance analysis and
    design aids, modeling of computer architecture
\end{IEEEkeywords}

\section{Introduction}
\subsection{Context}

When designing or modeling computer architectures, memory hierarchy
configurations are responsible for much of a system's performance across a
variety of workloads. Modern computer architectures face a fundamental problem,
known as a \textit{memory wall} \cite{wulf:1995}, the growing disparity between
the speed of \textsc{cpu} cores and main memory latencies. To mitigate this,
most if not all serious implementations of modern instruction set architectures
(\textsc{isa}s) employ complex memory hierarchies, split across multiple
levels, and which exploit locality. This hierarchy encompasses each core's
register file, different cache levels, as well as main and secondary memories,
and is one of the cornerstones of contemporary architecture design. Its
configuration parameters determine not only the maximum throughput of
high-performance systems, but also the energy efficiency and the viability of
low-power systems, as in single-board computers (\textsc{sbc}s) and distributed
systems.

The efficacy of a given memory hierarchy hinges on the ability to exploit
temporal and spatial localities in data accesses. However, the precise
configuration of memory hierarchy parameters---be they cache sizes,
associativity, latencies, replacement policies, block sizes, \textit{etc}.---to
achieve optimal throughput for a given workload remains a complex and
domain-dependent question, and requires detailed empirical exploration.

\subsection{Motivation}

Although much research has been done on the impact of cache configuration
parameters on a variety of workloads, most analyses measure only specific
individual metrics, focus on a small domain of specific workloads, or evaluate
a limited set of configuration parameters in isolation. Thus, it is often
difficult to systematically compare the relative effects of different cache
design choices on a wide array of workloads in a controlled, reproducible
manner. The lack of such multi-parameter evaluations motivates the need for a
methodology that enables consistent and interpretable explorations of the
design space.

\subsection{Objective}

In this work, we aim to conduct controlled architectural experiments using
select kernels from the PolyBench benchmark suite \cite{polybench:2016} in
order to investigate the impact of different cache memory configurations on
select performance metrics. A set of experimental scenarios is designed by
sweeping key parameters of the memory hierarchy, including cache size, cache
line size, and associativity levels for all cache levels. The experiments will
be conducted using the gem5 architectural simulator \cite{gem5-20:2007}. The
evaluation is based on the analysis of select performance metrics such as
multi-level miss rates and \textsc{mpki} (Misses Per Kilo Instructions) counts,
speedup across different simulated architectures, and memory-related stall
cycles. Through this analysis, we seek to provide a deeper understanding of how
variations in cache parameters effect changes in standard microarchitectural
metrics in the execution behavior of PolyBench benchmarks, and thus enable the
identification of potential optimization opportunities within the memory
hierarchy.

\subsection{Contribution}

We provide a systematic and reproducible evaluation of cache hierarchy
parameters across an extensible set of benchmark kernels, using the gem5
simulator. The main contributions are threefold. First, we perform a controlled
multi-parameter exploration of cache line size, associativity, and capacity,
enabling direct comparison of their relative impact on key microarchitectural
metrics such as miss rates, \textsc{mpki}, and \textsc{ipc}/\textsc{cpi}.
Second, we develop an automated simulation and data collection framework that
ensures consistency and reproducibility across many experimental
configurations. Third, we provide a cross-workload analysis that highlights how
different access patterns and dataset sizes influence the sensitivity of
performance to cache design choices, offering practical insights into the
trade-offs involved in memory hierarchy configuration.

\section{Related Work}

From an architectural optimization perspective, Khoshavi and DeMara
\cite{khoshavi:2018} investigate the use of emerging non-volatile memories
within the cache hierarchy. The authors propose the \textsc{rrap} strategy
(Read Reference Activity Persistent), a hybrid architecture that replaces
conventional \textsc{sram} with a combination of low- and high-retention
\textsc{stt}-\textsc{ram} (Spin-Transfer Torque Random Access Memory). Their
results indicate a 51.7\% reduction in L2 cache read miss rate, along with
improvements in area and energy consumption, without degrading system
performance.

From a methodological standpoint, Bueno et~al.\@\cite{bueno:2024} emphasize the
importance of experimental rigor and reproducibility in memory architecture
evaluation. Their work proposes a weight redefinition methodology for
simulation intervals based on memory activity (\textsc{mpki}). The authors
demonstrate that this approach reduces discrepancies in the analysis of
replacement policies by more than 30\%, highlighting that the reliability of
results strongly depends on the quality of the experimental procedure.

The validation of proposals in this domain relies on well-established
simulation infrastructures. The gem5 architectural simulator
\cite{gem5-20:2007} has become a fundamental tool for detailed architectural
modeling, enabling systematic exploration of design parameters. In parallel,
the investigation of spatial and temporal locality requires representative
workloads, for which the PolyBench suite \cite{polybench:2016} is widely
adopted due to its robust collection of mathematical kernels. The versatility
of gem5 has further driven recent research focused on Design Space Exploration
(\textsc{dse}), of which a representative example is the work of Shahid
et~al.\@\cite{shahid:2025}, in which the authors employ the framework to
improve the performance of the \textsc{risc}-V \textsc{cva6} core. Through
parametric simulations, they identify configurations capable of balancing
hardware cost and processor performance. This approach reinforces the
suitability of gem5 not only as a tool for evaluating traditional cache designs
but also as a key instrument for the rigorous validation of emerging and
extensible architectures.

\section{Methodology}
\subsection{Scope}

The goal is to evaluate cache hierarchy metrics of representative compute
kernels using the gem5 architectural simulator\cite{gem5-20:2007}. We establish
a three-level baseline cache configuration and simulate kernel executions while
sweeping over orthogonal parameters one at a time. The first of these
parameters is the cache hierarchy size, which is split into 11 presets (see
Table~\ref{tbl:cachesizes}) roughly inspired by real commercial
\textsc{cpu}s,\footnote{For the Intel Atom and IBM POWER10 configurations,
cache sizes were adjusted to the nearest power-of-two values to ensure
compatibility with the gem5 simulator. Apple Silicon cache sizes are taken from
the energy-efficiency cores.} where the Intel® Core™ i7-6700K was chosen as the
baseline. Next, we sweep the size of the cache lines over four power-of-two
values, from 32 to 256 bytes, where the baseline is 64 bytes; and finally we
sweep  L1, L2, and L3 associativities individually over 1-way (directly
mapped), 2-way, 4-way, 8-way, and 16-way configurations,\footnote{A 12-way
associativity configuration would be desirable, as it exists in the Zen~5
architecture, but gem5 expects power-of-two associativities.} where the
baselines were 8-way for L1 and L2, and 16-way for L3.

\begin{table}[!t]
	\centering
	\caption{Multi-level cache size presets}
	\label{tbl:cachesizes}
	\begin{tabular}{|l|c|c|c|c|c|}
		\hline
		\textbf{Size Preset}         & \textbf{L1I} & \textbf{L1D} & \textbf{L2} & \textbf{L3} \\
		\hline
		\hline
		i7-6700K \textit{(baseline)} & 32 KiB       & 32 KiB       & 256 KiB     & 8 MiB       \\
		Small Embedded               & 16 KiB       & 16 KiB       & 128 KiB     & 1 MiB       \\
		Cortex A78                   & 32 KiB       & 32 KiB       & 512 KiB     & 4 MiB       \\
		Intel Atom                   & 32 KiB       & 32 KiB       & 1 MiB       & 4 MiB       \\
		IBM POWER10                  & 32 KiB       & 32 KiB       & 2 MiB       & 8 MiB       \\
		Apple M1                     & 64 KiB       & 64 KiB       & 4 MiB       & 8 MiB       \\
		i9-9900K                     & 32 KiB       & 32 KiB       & 256 KiB     & 16 MiB      \\
		Ryzen 5600X                  & 32 KiB       & 32 KiB       & 512 KiB     & 32 MiB      \\
		Apple M2                     & 64 KiB       & 64 KiB       & 4 MiB       & 16 MiB      \\
		Ryzen 7700X                  & 32 KiB       & 32 KiB       & 1 MiB       & 32 MiB      \\
		Large Server                 & 64 KiB       & 64 KiB       & 2 MiB       & 64 MiB      \\
		\hline
	\end{tabular}
\end{table}

The choice of gem5 over other simulators, such as Sniper, is motivated by its
detailed and flexible modeling capabilities, as well as its broad support for
multiple \textsc{isa}s. While Sniper prioritizes high-speed simulation through
approximate interval-based models, gem5 enables cycle-accurate simulation and
fine-grained control over microarchitectural components. This level of detail
is helpful in isolating the effects of the cache parameters we seek to
investigatem, and in collecting detailed metrics.

The simulated gem5 system was configured to use a \textit{TimingSimpleCPU} with
the x86 \textsc{isa}, an in-order model with memory access timing, chosen
because the effects of out-of-order execution or prefetching on cache metrics
is not in the scope of this experiment. The system operates at a clock
frequency of 4~GHz and models a timing-based memory system over a physical
address space of 8~GiB. The processor is connected to a multi-level cache
hierarchy consisting of private L1 instruction and data caches, followed by
shared L2 and L3 caches. Each cache level is equipped with a finite number of
Miss Status Holding Registers (\textsc{mshr}s), to enable multiple outstanding
memory requests and allow limited memory-level parallelism. Communication
between cache levels and main memory is modeled using a hierarchical crossbar
interconnect, with dedicated crossbars between each level and a data width of
32 bytes. See Fig.~\ref{fig:system-diagram} for the system diagram generated by
the simulator.

\begin{figure*}[p]
	\centering
	\includegraphics[width=\textwidth]{../results/atax_baseline/config.dot.pdf}
	\caption{Diagram of simulated system configuration, automatically generated
	by gem5 when running a simulation.}
	\label{fig:system-diagram}
\end{figure*}

\subsection{Workload Selection}

We use five PolyBench Version~4.2.1~\cite{polybench:2016} kernels:
\textsc{ata}x, Floyd-Warshall, \textsc{gemm}, Jacobi-2D, and Seidel-2D. These
kernels capture a range of memory access patterns, including dense linear
algebra with regular streaming accesses (\textsc{ata}x, \textsc{gemm}), dynamic
programming with structured temporal reuse (Floyd-Warshall), and iterative
stencil computations with localized spatial reuse (Jacobi-2D, Seidel-2D). This
selection is intended to provide diverse locality characteristics while
maintaining manageable simulation time.

\subsubsection{\textsc{gemm} (General Matrix Multiply)} The \textsc{gemm}
kernel implements dense matrix multiplication, a core operation in
\textsc{blas} (Basic Linear Algebra Subprograms)~\cite{lawson:1979}, computing
a linear combination of matrix products:

\[
C_{ij} \leftarrow \beta C_{ij} + \alpha \sum_{k=0}^{N-1} A_{ik} B_{kj}
\]

From PolyBench~\cite{polybench:2016}:
\[
\text{Time: } O(n^3), \quad \text{Memory: } O(n^2)
\]

This kernel exhibits regular, dense access patterns with high spatial locality
and significant temporal reuse, particularly in the accumulation of $C_{ij}$.

\subsubsection{\textsc{ata}x (Matrix-Vector Transpose Product)} The
\textsc{ata}x kernel computes a matrix-vector product followed by a transpose
multiplication, a common pattern in iterative solvers.

\[
y = A^T (A x)
\]

From PolyBench~\cite{polybench:2016}:
\[
\text{Time: } O(n^2), \quad \text{Memory: } O(n^2)
\]

The kernel combines row-major streaming accesses in the forward multiplication
with column-wise accesses in the transpose phase, which results in mixed
spatial locality and limited temporal reuse.

\subsubsection{Floyd-Warshall (All-Pairs Shortest Paths)} The Floyd-Warshall
algorithm computes shortest paths between all pairs of vertices in a weighted
graph using dynamic programming.\cite{floyd:1962}\cite{warshall:1962}

\[
D_{ij}^{(k)} = \min\left(D_{ij}^{(k-1)},\; D_{ik}^{(k-1)} + D_{kj}^{(k-1)}\right)
\]

From PolyBench~\cite{polybench:2016}:
\[
\text{Time: } O(n^3), \quad \text{Memory: } O(n^2)
\]

The algorithm performs dense, regular accesses over a 2D matrix with strong
temporal reuse across iterations of the outer loop.

\subsubsection{Jacobi-2D (Iterative Stencil Method)} The Jacobi-2D kernel
applies an iterative stencil computation over a 2D grid, commonly used in
numerical solutions of partial differential equations.\cite{golub:2013}

\[
A_{i,j}^{(t+1)} =
\frac{1}{5}\left(
A_{i,j}^{(t)} +
A_{i-1,j}^{(t)} +
A_{i+1,j}^{(t)} +
A_{i,j-1}^{(t)} +
A_{i,j+1}^{(t)}
\right)
\]

From PolyBench~\cite{polybench:2016}:
\[
\text{Time: } O(n^3), \quad \text{Memory: } O(n^2)
\]

The kernel exhibits structured spatial locality due to stencil-based neighbor
accesses, with moderate temporal locality across time steps.

\subsubsection{Seidel-2D (Gauss--Seidel Stencil)} The Seidel-2D kernel
implements an in-place stencil computation using a Gauss-Seidel update scheme,
where newly computed values are immediately reused within the same iteration.

\[
A_{i,j} \leftarrow \frac{1}{9}
\sum_{p=-1}^{1} \sum_{q=-1}^{1} A_{i+p,j+q}
\]

From PolyBench~\cite{polybench:2016}:
\[
\text{Time: } O(n^3), \quad \text{Memory: } O(n^2)
\]

Compared to Jacobi-2D, the in-place update introduces loop-carried
dependencies, increasing temporal locality but reducing
parallelism.\cite{golub:2013}

\subsection{Collection of Metrics}

Each workload was compiled as a statically linked 64-bit \textsc{elf}
executable targeting the x86-64 \textsc{isa} and the Linux system call
interface. The musl C standard library~\cite{musl} was used to ensure minimal
and deterministic runtime dependencies. All binaries were linked against gem5's
\textit{libm5} library. Instrumentation points were inserted via the
\textit{libm5} \textsc{api} to delimit the region of interest (\textsc{roi}).
Specifically, gem5 annotation calls were placed immediately before and after
the kernel invocation, which ensures precise control over when statistics
collection begins and ends.

For each run we record a standard set of counters from gem5's generated
\texttt{stats.txt}: processor throughput (\textsc{ipc}/\textsc{cpi}); per-level
cache hits and misses and their miss rates (L1I, L1D, L2, L3); and two derived
quantities: \textsc{mpki} (misses per kilo instructions) for each level, and
the total miss-latency (sum of overall miss latencies across levels) used as a
proxy for memory stalls. For comparisons, we use a consistent per-workload
baseline preset: categorical presets are shown as deltas or speedup
(\(\mathrm{\textsc{ipc}}/\mathrm{\textsc{ipc}}_{\text{baseline}}\)), while
numeric sweeps (e.g., cache line size) are plotted as absolute values, with
logarithmic scales applied where appropriate to span orders of magnitude.

\subsection{Datasets}\label{sub:dataset-size}

PolyBench kernels are parameterized by dataset size through preprocessor
definitions (e.g., \texttt{SMALL\_\-DATA\-SET}, \texttt{MEDIUM\_\-DATA\-SET},
\texttt{LARGE\_\-DATA\-SET}), with \texttt{LARGE\_\-DATA\-SET} selected by
default if none is specified. In this work, specifically, dataset sizes were
explicitly configured per kernel in order to balance simulation fidelity with
tractable execution time under detailed timing simulation.
Table~\ref{tab:datasets} summarizes the selected dataset for each workload
along with the corresponding problem dimensions.

\begin{table}[!t]
	\centering
	\caption{Dataset configurations used for each workload.}
	\begin{tabular}{|l|c|c|}
		\hline
		\textbf{Kernel} & \textbf{Dataset} & \textbf{Dimensions} \\
		\hline
		\hline
		\textsc{ata}x   & LARGE            & $M = 1900$, $N = 2100$ \\
		\textsc{gemm}   & MEDIUM           & $N_I = 200$, $N_J = 220$, $N_K = 240$ \\
		Floyd-Warshall  & SMALL            & $N = 180$ \\
		Jacobi-2D       & MEDIUM           & $N = 250$, $T = 100$ \\
		Seidel-2D       & MEDIUM           & $N = 400$, $T = 100$ \\
		\hline
	\end{tabular}
	\label{tab:datasets}
\end{table}

Preliminary simulations using the baseline cache configuration indicated, upon
register inspection after partial simulation progress, that several kernels
with the default \texttt{LARGE\_DATASET} would require prohibitively long
execution times under cycle-accurate simulation. Dataset sizes were therefore
reduced iteratively until each workload could complete within a practical time
budget (on the order of days). In particular, Floyd-Warshall exhibits cubic
time complexity with a large constant factor, which necessitated the use of the
\texttt{SMALL\_DATASET} configuration.

\subsection{Reproducibility and Build Pipeline}

We express the end-to-end workflow as an ordered, idempotent build graph in the
Zig Build System (tied to Zig version~0.15.2~\cite{zig:2024}). The graph
encodes, in dependency order: dependency checks; a pinned Python runtime via
\texttt{uv} and a project virtual environment; gem5 submodule initialization
and compilation; the \textit{libm5} library build; workload compilation and
linking; the complete simulation sweep; figure generation; and the final
\LaTeX{} report build.

Versions are pinned to minimize drift. Python is fixed to 3.14.3 via
\texttt{uv}, and all Python packages are frozen in \texttt{requirements.txt}.
The gem5 source is vendored as a Git submodule and built at the exact recorded
commit (\texttt{7a2b0e4}). Workloads are statically linked x86-64 \textsc{elf}
binaries against musl to reduce host-library variance. The C/C++ toolchain is
provided by Zig's bundled Clang (\texttt{zig~cc}/\texttt{zig~c++}). All
invocations of gem5, simulations, and analysis are routed through the project
\texttt{venv} to ensure a consistent runtime across machines.

Simulations are dispatched by a single Python orchestration script that
enumerates the parameter grid, manages a worker pool, can pin workers to
physical CPU cores to reduce noise, and writes each simulation's collected
metrics under deterministic paths. The runner uses completion markers and
key-counter checks to make the sweep idempotent and resumable. Together, the
build graph and the runner provide a concise reproduction path. All source code
and artifacts are publicly available in a versioned
repository~\cite{repo:2026}.

\section{Results}
\subsection{Overview}

Overall, only a small subset of the values explored in our parameter sweep had
a significant impact on the sampled metrics. The most dominant parameter across
all workloads by far was the cache line size, wherein doubling the line size
tends to significantly reduce L1D miss rate, often approaching a halving for
highly sequential access patterns (\textsc{ipc} improves accordingly), because
larger lines increase spatial locality utilization per miss, and associativity
effects are mostly tied to cache entry conflicts within a set. As such,
associativity effects were largely flat, and the capacity presets impact
performance selectively, and only when dataset footprints\footnote{To be
understood as the kernel's hot working set.} fit in the larger caches.

\begin{figure}[!t]
    \centering
    \caption{\textsc{ipc} vs. cache line sizes across kernels}
    \includegraphics[width=\linewidth]{../figures/comparison_ipc_vs_cache_line.pdf}
\end{figure}

\begin{figure}[!t]
    \centering
    \caption{L1D miss rate vs. cache line sizes across kernels}
    \includegraphics[width=\linewidth]{../figures/comparison_l1d_miss_rate_vs_cache_line.pdf}
\end{figure}

\subsection{Cache Line Size Effects}\label{sub:cacheline}

Larger lines better exploit spatial locality by coalescing contiguous words per
tag. Thus, fewer L1D misses propagate as lower miss latencies and lead to
\textsc{cpi} gains. To point out some examples,
Fig.~\ref{fig:atax-speedup-line} shows that, for the \textsc{ata}x kernel,
contiguous accesses give high spatial locality, and doubling line size captures
more contiguous work per refill, which in turn reduces the number of refills.
On Fig.~\ref{fig:jacobi-mpki-line}, Jacobi-2D shows the same spatial-locality
response, and on Fig.~\ref{fig:seidel-stalls-line}, we see that Seidel-2D sees
dramatic reductions in miss latencies with larger lines.\footnote{In these
latter two, it's worth noting that \textsc{cpi} gains are present but smaller,
because a significant fraction of execution is compute. A larger dataset would
probably correspond to a more dramatic \textsc{cpi} reduction.} Finally,
Fig.~\ref{fig:gemm-cpi-line} shows modest \textsc{cpi} improvements for
\textsc{gemm}, due to inherent temporal locality and compute intensity.

\begin{figure}[!t]
    \centering
	\caption{\textbf{\textsc{ata}x} relative speedup vs. cache line sizes.
	Correlation indicates a strongly spatial, memory-sensitive kernel.}
    \label{fig:atax-speedup-line}
    \includegraphics[width=\linewidth]{../figures/atax_speedup_vs_cache_line.pdf}
\end{figure}

\begin{figure}[!t]
    \centering
    \caption{\textbf{Jacobi-2D} L3 \textsc{mpki} halve for each doubling of
    cache line size.}
    \label{fig:jacobi-mpki-line}
    \includegraphics[width=\linewidth]{../figures/jacobi-2d_mpki_vs_cache_line.pdf}
\end{figure}

\begin{figure}[!t]
    \centering
    \caption{\textbf{Seidel-2D} latency also decreases significantly for each
    doubling of cache line size.}
    \label{fig:seidel-stalls-line}
    \includegraphics[width=\linewidth]{../figures/seidel-2d_memory_stalls_vs_cache_line.pdf}
\end{figure}

\begin{figure}[!t]
    \centering
	\caption{\textbf{\textsc{gemm}} \textsc{cpi} shows gains are modest,
	because kernel is compute-heavy with non-zero memory sensitivity.}
    \label{fig:gemm-cpi-line}
    \includegraphics[width=\linewidth]{../figures/gemm_cpi_vs_cache_line.pdf}
\end{figure}

\subsection{Associativity Changes Least Significant}

Cache associativity was by far the least significant cache parameter, and the
impact was almost negligible (a few tenths of a percent). Mechanically,
associativity reduces conflict misses by allowing multiple cache lines that map
to the same set to coexist. In our kernels, most accesses are regular and
spread across many sets. That is because strides are low or simple stencils are
used. Thus, conflict pressure is low, suggesting that misses are dominated by
compulsory and capacity effects. The only exception is when associativity is
one (the cache is directly mapped), due to the fact that \textit{all} conflicts
will lead to an eviction. This can be seen in Fig.~\ref{fig:comp-ipc-l1} as
slightly lower instructions per cycle when the L1 caches are directly mapped.

\begin{figure}[!t]
    \centering
	\caption{\textsc{ipc} vs. L1 associativity across kernels shows negligible
	impact on performance.}
    \label{fig:comp-ipc-l1}
    \includegraphics[width=\linewidth]{../figures/comparison_ipc_vs_l1_assoc.pdf}
\end{figure}

\subsection{Capacity Presets}

The presets change the nominal capacity of L1, L2, and L3. Larger overall
capacities reduce misses \textit{when} the kernel's hot working set is small
enough that the last-level cache (\textsc{llc}) can hold a significant fraction
of the set. When the set is significantly larger than the \textsc{llc} can
hold, or when most misses are compulsory, such that spatial locality
optimizations (e.g., larger lines) are more impactful, as we exemplified in
\ref{sub:cacheline}, capacity changes have little effect.
Fig.~\ref{fig:comp-ipc-size} shows which workloads had most gains, and
Fig.~\ref{fig:atax-speed-size} shows the effect on the \textsc{ata}x kernel,
which had the most dramatic improvement.

\begin{figure}[!t]
    \centering
	\caption{$\Delta$\textsc{ipc} across cache configurations, showing that
	only a subset produces noticeable gains.}
    \label{fig:comp-ipc-size}
    \includegraphics[width=\linewidth]{../figures/comparison_ipc_vs_cache_config.pdf}
\end{figure}

\begin{figure}[!t]
    \centering
	\caption{\textbf{\textsc{ata}x} speedup as total cache size grows was the
	highest among all tested kernels.}
    \label{fig:atax-speed-size}
    \includegraphics[width=\linewidth]{../figures/atax_speedup_vs_cache_config.pdf}
\end{figure}

\subsection{Some Observations}

\subsubsection{Instruction cache misses} Across all 155 simulations, the total
miss count for the instruction cache was \textit{always} zero. This is to be
expected, because by the time the workload binary calls the \textit{libm5}
\textsc{api} to start counting statistics, the working set of instructions has
already been fetched into the instruction cache and easily fits within its
capacity: the relocatable object corresponding to the main translation unit of
each workload is on the order of a few kilobytes.

\subsubsection{Impact of dataset sizes} As we detailed in
\ref{sub:dataset-size}, different workloads were set to use different dataset
size configurations due to time constraints. The \textsc{ata}x kernel was
compiled with the \texttt{LARGE\_DATASET} definition, the Floyd-Warshall kernel
with \texttt{SMALL\_DATASET}, and the remaining kernels with
\texttt{MEDIUM\_DATASET}. This may explain why the \textsc{ata}x simulations
exhibited the largest absolute jumps in \textsc{mpki}, miss rates, total
latencies, and \textsc{cpi}, whereas Floyd-Warshall exhibited the smallest
jumps: larger dataset sizes may increase working-set pressure, such that, if
the working set is near the cache size boundary, changes in total cache
capacity and line size may materially change miss behavior.

\section{Conclusion}

We evaluated the sensitivity of cache hierarchies across five PolyBench
kernels under a controlled and reproducible simulation methodology.
By decoupling realistic capacity presets from orthogonal cache parameters,
we observed that:
\begin{enumerate*}[label=(\alph*), itemjoin={;\ }, itemjoin*={; and\ }]
    \item cache line size is the dominant factor, consistently reducing L1D
    miss rates and improving \textsc{cpi} in memory-sensitive regimes
    \item associativity has limited impact for regular access patterns, with
    noticeable effects primarily in direct-mapped configurations
    \item increased capacity provides benefits selectively, mainly when the
    working set approaches the \textsc{llc} capacity.
\end{enumerate*}
These results suggest that, for regular workloads with predictable access
patterns, spatial locality optimizations yield the largest performance gains,
while conflict-related effects are comparatively minor.

The interpretation of these results is subject to several limitations. First,
all experiments were conducted using an in-order core model, which does not
capture the effects of out-of-order execution or memory-level parallelism.
Second, dataset sizes were adjusted to ensure tractable simulation times, which
may reduce pressure on the memory hierarchy relative to real-world scenarios.
Finally, the selected workloads exhibit regular access patterns, and therefore
may not reflect the behavior of irregular or indirection-intensive
applications.

Reproducibility is a central aspect of this work. The entire experimental
pipeline---including dependency management, workload compilation, simulation
orchestration, and analysis---is encoded in a deterministic build graph with
pinned versions, enabling consistent regeneration of results across
environments. Future works may consider scaling datasets to stress
\textsc{llc}s and main memory more aggressively, extending the workload set to
include irregular applications, exploring additional cache design dimensions
such as replacement and prefetching policies, and evaluating the same parameter
space on out-of-order cores.

\bibliographystyle{IEEEtran}
\bibliography{references}

\end{document}
